\documentclass[../../main.tex]{subfiles}

\begin{document}

\section{一致有界原理}\label{section:ch13-一致有界}

作为 Baire 范畴定理的推论, 我们证明: 若完备度量空间上的连续函数族是逐点有界的, 则存在开集使得该族函数在其上是一致有界的. 这有以下富有成效的关于线性算子族的推论.

\begin{theorem}[一致有界原理]\label{theorem:ch13-一致有界}
令 $X$ 为 Banach 空间而 $Y$ 是赋范线性空间. 考虑算子族 $\mathcal{F}\subseteq\mathcal{L}(X,Y)$. 假定 $\mathcal{F}$ 逐点有界, 即对 $X$ 中的每个 $x$, 存在常数 $M_{x}\geqslant 0$ 使得
\begin{align*}
\text{对所有 }T\in\mathcal{F},\quad\|T(x)\|\leqslant M_{x}.
\end{align*}
则 $\mathcal{F}$ 在以下意义下是一致有界的: 存在常数 $M\geqslant 0$, 使得对所有 $\mathcal{F}$ 中的 $T$, $\|T\|\leqslant M$.
\end{theorem}

\begin{proof}
对每个 $T\in\mathcal{F}$, 定义为 $f_{T}(x)=\|T(x)\|$ 的实值函数 $f_{T}:X\to\mathbf{R}$ 是 $X$ 上的一个实值连续函数. 由于该连续函数族在 $X$ 上是逐点有界的而度量空间 $X$ 是完备的, 根据第 10 章的定理 6, 存在 $X$ 中的开球 $B(x_{0},r)$ 与常数 $C\geqslant 0$, 使得
\begin{align*}
\text{对所有 }x\in B(x_{0},r)\text{ 与 }T\in\mathcal{F},\quad\|T(x)\|\leqslant C.
\end{align*}
因此, 对每个 $T\in\mathcal{F}$,
\begin{align*}
\text{对所有 }x\in B(0,r),\quad\|T(x)\|=\|T([x+x_{0}]-x_{0})\|\\
\leqslant\|T(x+x_{0})\|+\|T(x_{0})\|\leqslant C+M_{x_{0}}.
\end{align*}
因此, 设 $M=(1/r)(C+M_{x_{0}})$, 对 $\mathcal{F}$ 中的所有 $T$ 我们有 $\|T\|\leqslant M$.
\end{proof}

\begin{theorem}[Banach-Saks-Steinhaus 定理]\label{theorem:ch13-Banach-Steinhaus}
令 $X$ 为 Banach 空间, $Y$ 是赋范线性空间, 而 $\{T_{n}:X\to Y\}$ 是连续线性算子序列. 假定对每个 $x\in X$,
\begin{align}
\lim_{n\to\infty}T_{n}(x)\text{ 在 }Y\text{ 中存在}.\label{align:ch13-逐点极限存在}
\end{align}
则算子序列 $\{T_{n}:X\to Y\}$ 是一致有界的. 此外, 对所有 $x\in X$, 定义为
\begin{align*}
T(x)=\lim_{n\to\infty}T_{n}(x)
\end{align*}
的算子 $T:X\to Y$ 是线性、连续的, 且
\begin{align*}
\|T\|\leqslant\liminf\|T_{n}\|.
\end{align*}
\end{theorem}

\begin{proof}
线性算子的逐点极限是线性的. 因此 $T$ 是线性的. 我们从一致有界原理推出序列 $\{T_{n}\}$ 是一致有界的. 因此 $\liminf\|T_{n}\|$ 是有限的. 令 $x$ 属于 $X$. 根据 $Y$ 上的范数的连续性, $\lim_{n\to\infty}\|T_{n}(x)\|\to\|T(x)\|$. 由于对所有 $n$, $\|T_{n}(x)\|\leqslant\|T_{n}\|\cdot\|x\|$, 我们也有 $\|T(x)\|\leqslant\liminf\|T_{n}\|\cdot\|x\|$. 因此 $T$ 是有界的且 $\|T\|\leqslant\liminf\|T_{n}\|$. 有界线性算子是连续的.
\end{proof}

在 $Y$ 是 Banach 空间的情形, \eqref{align:ch13-逐点极限存在}等价于以下断言: 对每个 $x\in X$, $\{T_{n}(x)\}$ 是 $Y$ 中的 Cauchy 序列.

\begin{exercise}\label{exercise:ch13-38}
作为 Baire 范畴定理的一个推论, 我们证明作为完备度量空间上的连续映射序列的逐点极限的实值映射, 必须在它的定义域的某个稠密子集的所有点连续. 调整上面的证明, 使得它适用于映入任何度量空间的映射. 用此证明 Banach 空间上的连续线性算子序列的逐点极限有一个极限, 它在某个点是连续的, 因此根据线性它是连续的.
\end{exercise}

\begin{exercise}\label{exercise:ch13-39}
令 $\{f_{n}\}$ 为 $L^{\infty}[a,b]$ 中的序列. 假定对每个 $g\in L^{1}[a,b]$, $\lim_{n\to\infty}\int_{a}^{b}g\cdot f_{n}$ 存在. 证明存在函数 $f\in L^{\infty}[a,b]$, 使得对所有 $g\in L^{1}[a,b]$, $\lim_{n\to\infty}\int_{a}^{b}g\cdot f_{n}=\int_{a}^{b}g\cdot f$.
\end{exercise}

\begin{exercise}\label{exercise:ch13-40}
令 $X$ 为定义在 $\mathbf{R}$ 上的所有多项式组成的线性空间. 对 $p\in X$, 定义 $\|p\|$ 为 $p$ 的系数的绝对值的和. 证明这是 $X$ 上的范数. 对每个 $n$, 定义 $\psi_{n}:X\to\mathbf{R}$ 为 $\psi_{n}(p)=p^{(n)}(0)$. 用 $\mathcal{L}(X,\mathbf{R})$ 中的序列 $\{\psi_{n}\}$ 的性质证明 $X$ 不是 Banach 空间.
\end{exercise}

\end{document}